\begin{pastexam}[label={px:ccu-108-swa-09}]{中正資工所甲組 108 學年度 第 10 題}
（問答）

(1\%) What is the name of such representation when a directed graph with edges ($u\to v$) is
defined as follows?

\begin{center}
\begin{tabular}{|l|l|l|l|}\hline
$v \backslash u$ & 1 & 2 & 3\\\hline
1 & 0 & 0 & 0\\\hline
2 & 20 & 0 & -30\\\hline
3 & 40 & 0 & 0\\\hline
\end{tabular}
\end{center}
\end{pastexam}

\begin{pastsoln}
答案：（加權的）鄰接矩陣（adjacency matrix），也有書叫成本矩陣（cost matrix）。

$n$ 個頂點的圖用一個 $n \times n$ 的矩陣表示，格子存那條邊的權重，沒有邊的格子存一個代表「沒有邊」
的值，這就是加權圖的鄰接矩陣（資結書第~21.2~節）。 讀這張表要注意兩件事：
\begin{itemize}
\item 左上角寫 $v \backslash u$：\textbf{列}是終點 $v$、\textbf{行}是起點 $u$，格子 $(v, u)$ 是
$u \to v$ 的權重。 這和常見的「第 $u$ 列第 $v$ 行存 $u \to v$」正好轉置。 所以第 $2$ 列第 $1$ 行的
$20$ 是邊 $1 \to 2$，不是 $2 \to 1$。
\item 這張表用 $0$ 表示沒有邊（對角線也是 $0$）。 所以圖只有三條邊：$1 \to 2$（$20$）、
$1 \to 3$（$40$）、$3 \to 2$（$-30$）。 其中有一條負邊，這是下一題（第 11 題）的伏筆。
\end{itemize}
\end{pastsoln}
